% GENERATED FILE. Edit canonical lessons, then run npm run generate:books.
% canonical-source-hash: fnv1a64:a8513d0115574a8f
% canonical-lessons: IT-C402-rete, IT-C402-connessione, IT-C402-connettersi, IT-C402-navigare, IT-C402-link, IT-R402-ripasso-395-402

\chapter{In rete}
\label{ch:it-in-rete}
\hlchaptermodality{402}

\begin{chapteropening}
\textbf{By the end of this chapter:} \emph{I can talk about being online: the network and the connection, connecting and browsing, and a link.}

Complete the last lesson of chapter 402: I can talk about being online: the network and the connection, connecting and browsing, and a link.
\end{chapteropening}

\section[la rete]{la rete — the network; the internet; the net}

\label{lesson:IT-C402-rete}



\begin{quote}
Before the new one: say the Italian for to paste, then the Italian for the backup.
\end{quote}

\subsection*{You'll want to know: la rete}
\textbf{la rete} — \textquotedblleft{}the network; the internet; the net\textquotedblright{}.

\begin{quote}
Ho trovato la ricetta in rete.
\end{quote}

\begin{quote}
\emph{I found the recipe online.}
\end{quote}

\emph{In rete} — online. In football, \emph{la rete} is the goal net, and a goal.

\begin{cousinweb}
\emph{rete} is inherited from Latin \emph{rēte} \textquotedblleft{}a net\textquotedblright{} (Lewis \& Short): a fishing net, then any network of threads, roads or computers.
\end{cousinweb}

\begin{grammarlens}[title={What you've built}]
One more word for telling, explaining and describing at B1.
\end{grammarlens}

\subsection*{Your turn}
\begin{itemize}
\raggedright
  \item \emph{Say it:} \emph{la rete}
  \item \emph{Say it:} \emph{la rete}, once more
  \item \emph{Say it:} a sentence of your own with \emph{la rete}
  \item \emph{Recall:} say the Italian for to paste, then the Italian for the backup, then say \emph{la rete} again
\end{itemize}

\begin{tcolorbox}[breakable,colback=teal!4,colframe=teal!35!black,title={Before you move on}]
What does la rete mean? (\textquotedblleft{}The network; the internet; the net\textquotedblright{}.) Say it once more.
\end{tcolorbox}
\section[la connessione]{la connessione — the connection}

\label{lesson:IT-C402-connessione}



\begin{quote}
Before the new one: say the Italian for the backup, then the Italian for the network.
\end{quote}

\subsection*{You'll want to know: la connessione}
\textbf{la connessione} — \textquotedblleft{}the connection\textquotedblright{}.

\begin{quote}
Se la connessione è lenta, il video si ferma.
\end{quote}

\begin{quote}
\emph{If the connection is slow, the video stops.}
\end{quote}

Also a tie between ideas.

\begin{cousinweb}
\emph{connessione} is a learned borrowing from Latin \emph{cōnexiō, cōnexiōnis} \textquotedblleft{}a binding together\textquotedblright{} (EtymDB), from \emph{cōnectere} \textquotedblleft{}to tie together\textquotedblright{}, built on \emph{nectere} \textquotedblleft{}to tie\textquotedblright{}.
\end{cousinweb}

\begin{grammarlens}[title={What you've built}]
One more word for telling, explaining and describing at B1.
\end{grammarlens}

\subsection*{Your turn}
\begin{itemize}
\raggedright
  \item \emph{Say it:} \emph{la connessione}
  \item \emph{Say it:} \emph{la connessione}, once more
  \item \emph{Say it:} a sentence of your own with \emph{la connessione}
  \item \emph{Recall:} say the Italian for the backup, then the Italian for the network, then say \emph{la connessione} again
\end{itemize}

\begin{tcolorbox}[breakable,colback=teal!4,colframe=teal!35!black,title={Before you move on}]
What does la connessione mean? (\textquotedblleft{}The connection\textquotedblright{}.) Say it once more.
\end{tcolorbox}
\section[connettersi]{connettersi — to connect, to go online}

\label{lesson:IT-C402-connettersi}



\begin{quote}
Before the new one: say the Italian for the network, then the Italian for the connection.
\end{quote}

\subsection*{You'll want to know: connettersi}
\textbf{connettersi} — \textquotedblleft{}to connect, to go online\textquotedblright{}.

\begin{quote}
Se ti connetti al wifi dell'albergo, non paghi.
\end{quote}

\begin{quote}
\emph{If you connect to the hotel wifi, you do not pay.}
\end{quote}

\emph{Connettersi a} — to connect to. Past participle \emph{connesso}.

\begin{cousinweb}
\emph{connettersi} comes from Latin \emph{cōnectere} \textquotedblleft{}to tie together\textquotedblright{} (Lewis \& Short), built from \emph{con-} \textquotedblleft{}together\textquotedblright{} and \emph{nectere} \textquotedblleft{}to tie\textquotedblright{}.
\end{cousinweb}

\begin{grammarlens}[title={What you've built}]
One more word for telling, explaining and describing at B1.
\end{grammarlens}

\subsection*{Your turn}
\begin{itemize}
\raggedright
  \item \emph{Say it:} \emph{connettersi}
  \item \emph{Say it:} \emph{connettersi}, once more
  \item \emph{Say it:} a sentence of your own with \emph{connettersi}
  \item \emph{Recall:} say the Italian for the network, then the Italian for the connection, then say \emph{connettersi} again
\end{itemize}

\begin{tcolorbox}[breakable,colback=teal!4,colframe=teal!35!black,title={Before you move on}]
What does connettersi mean? (\textquotedblleft{}To connect, to go online\textquotedblright{}.) Say it once more.
\end{tcolorbox}
\section[navigare]{navigare — to browse, to surf; to sail}

\label{lesson:IT-C402-navigare}



\begin{quote}
Before the new one: say the Italian for the connection, then the Italian for to connect.
\end{quote}

\subsection*{You'll want to know: navigare}
\textbf{navigare} — \textquotedblleft{}to browse, to surf; to sail\textquotedblright{}.

\begin{quote}
Passo troppo tempo a navigare su internet.
\end{quote}

\begin{quote}
\emph{I spend too much time browsing the internet.}
\end{quote}

First to sail: \emph{navigare sul fiume}.

\begin{cousinweb}
\emph{navigare} is a learned borrowing from Latin \emph{nāvigāre} \textquotedblleft{}to sail\textquotedblright{} (Lewis \& Short), built on \emph{nāvis} \textquotedblleft{}a ship\textquotedblright{}.
\end{cousinweb}

\begin{grammarlens}[title={What you've built}]
One more word for telling, explaining and describing at B1.
\end{grammarlens}

\subsection*{Your turn}
\begin{itemize}
\raggedright
  \item \emph{Say it:} \emph{navigare}
  \item \emph{Say it:} \emph{navigare}, once more
  \item \emph{Say it:} a sentence of your own with \emph{navigare}
  \item \emph{Recall:} say the Italian for the connection, then the Italian for to connect, then say \emph{navigare} again
\end{itemize}

\begin{tcolorbox}[breakable,colback=teal!4,colframe=teal!35!black,title={Before you move on}]
What does navigare mean? (\textquotedblleft{}To browse, to surf; to sail\textquotedblright{}.) Say it once more.
\end{tcolorbox}
\section[il link]{il link — the link}

\label{lesson:IT-C402-link}



\begin{quote}
Before the new one: say the Italian for to connect, then the Italian for to browse.
\end{quote}

\subsection*{You'll want to know: il link}
\textbf{il link} — \textquotedblleft{}the link\textquotedblright{}.

\begin{quote}
Ti mando il link: clicca e guarda.
\end{quote}

\begin{quote}
\emph{I will send you the link: click and have a look.}
\end{quote}

Masculine, invariable: \emph{i link}. Also \emph{il collegamento}.

\begin{cousinweb}
\emph{link} was borrowed from English \emph{link} \textquotedblleft{}a ring of a chain\textquotedblright{} (EtymDB), a word usually traced to Old Norse.
\end{cousinweb}

\begin{grammarlens}[title={What you've built}]
One more word for telling, explaining and describing at B1.
\end{grammarlens}

\subsection*{Your turn}
\begin{itemize}
\raggedright
  \item \emph{Say it:} \emph{il link}
  \item \emph{Say it:} \emph{il link}, once more
  \item \emph{Say it:} a sentence of your own with \emph{il link}
  \item \emph{Recall:} say the Italian for to connect, then the Italian for to browse, then say \emph{il link} again
\end{itemize}

\begin{tcolorbox}[breakable,colback=teal!4,colframe=teal!35!black,title={Before you move on}]
What does il link mean? (\textquotedblleft{}The link\textquotedblright{}.) Say it once more.
\end{tcolorbox}
\section[(dialogue)]{Ripasso dei capitoli 395–402}

\label{lesson:IT-R402-ripasso-395-402}



\begin{quote}
Say the Italian for the saving and the investment.
\end{quote}

\subsection*{Your turn}
Say these aloud:

\begin{itemize}
\raggedright
  \item the saving, the investment, the interest, the rate, the inflation
  \item the cash, the balance, the cheque, the payment, the deposit
  \item the guarantee, to guarantee, the invoice, the quote, good value
  \item to earn, the earnings, the profit, the loss, to go bankrupt
  \item the online account, the password, the user, to log in, to sign up
  \item the app, to install, to uninstall, the device, the software
  \item the computer file, to delete, to copy, to paste, the backup
  \item the network, the connection, to connect, to browse, the link
\end{itemize}

\begin{tcolorbox}[breakable,colback=teal!4,colframe=teal!35!black,title={Before you move on}]
The saving? (\textbf{il risparmio}.) The investment? (\textbf{l'investimento}.) The interest? (\textbf{l'interesse}.) The rate? (\textbf{il tasso}.) The inflation? (\textbf{l'inflazione}.) The cash? (\textbf{i contanti}.) The balance? (\textbf{il saldo}.) The cheque? (\textbf{l'assegno}.) The payment? (\textbf{il versamento}.) The deposit? (\textbf{il deposito}.) The guarantee? (\textbf{la garanzia}.) To guarantee? (\textbf{garantire}.) The invoice? (\textbf{la fattura}.) The quote? (\textbf{il preventivo}.) Good value? (\textbf{conveniente}.) To earn? (\textbf{guadagnare}.) The earnings? (\textbf{il guadagno}.) The profit? (\textbf{il profitto}.) The loss? (\textbf{la perdita}.) To go bankrupt? (\textbf{fallire}.) The online account? (\textbf{l'account}.) The password? (\textbf{la password}.) The user? (\textbf{l'utente}.) To log in? (\textbf{accedere}.) To sign up? (\textbf{registrarsi}.) The app? (\textbf{l'applicazione}.) To install? (\textbf{installare}.) To uninstall? (\textbf{disinstallare}.) The device? (\textbf{il dispositivo}.) The software? (\textbf{il software}.) The computer file? (\textbf{il file}.) To delete? (\textbf{cancellare}.) To copy? (\textbf{copiare}.) To paste? (\textbf{incollare}.) The backup? (\textbf{il backup}.) The network? (\textbf{la rete}.) The connection? (\textbf{la connessione}.) To connect? (\textbf{connettersi}.) To browse? (\textbf{navigare}.) The link? (\textbf{il link}.)
\end{tcolorbox}
